\providecommand{\Res}[2]{\qty[round-mode=figures,round-precision=3]{\fpeval{#1}}{#2}}
\providecommand{\Val}[1]{\num[round-mode=figures,round-precision=4]{\fpeval{#1}}}
\providecommand{\ResSci}[2]{\qty[round-mode=figures,round-precision=3,exponent-mode=scientific]{\fpeval{#1}}{#2}}
\subsubsection*{Question a) Fréquence}
\begin{enumerate}[label=\arabic*)]
    \item \textbf{Variables connues et recherchées :} $\lambda = \qty{\v{lam}}{nm} = \qty{\v{lam}e-9}{m}$, $c = \qty{3.00e8}{m/s}$ ; variable recherchée : $f$
    \item \textbf{Équations pertinentes :}
    \[ c = \lambda f \]
    \item \textbf{Isolation de la variable recherchée :}
    \[ f = \frac{c}{\lambda} \]
    \item \textbf{Résolution numérique :}
    \[ f = \frac{\qty{3.00e8}{m/s}}{\qty{\v{lam}e-9}{m}} = \ResSci{3e8/\vsi{lam}}{Hz} \]
    \item \textbf{Réponse :} $f = \mathbf{\ResSci{3e8/\vsi{lam}}{Hz}}$.
\end{enumerate}

\subsubsection*{Question b) Période}
\begin{enumerate}[label=\arabic*)]
    \item \textbf{Variables connues et recherchées :} $f = \ResSci{3e8/\vsi{lam}}{Hz}$ ; variable recherchée : $T$
    \item \textbf{Équations pertinentes :}
    \[ T = \frac{1}{f} = \frac{\lambda}{c} \]
    \item \textbf{Isolation de la variable recherchée :} l'équation est déjà isolée.
    \item \textbf{Résolution numérique :}
    \[ T = \frac{\qty{\v{lam}e-9}{m}}{\qty{3.00e8}{m/s}} = \ResSci{\vsi{lam}/3e8}{s} = \Res{1e15*\vsi{lam}/3e8}{fs} \]
    \item \textbf{Réponse :} $T = \mathbf{\Res{1e15*\vsi{lam}/3e8}{fs}}$.
\end{enumerate}

\subsubsection*{Question c) Longueurs d'onde dans un millimètre}
\begin{enumerate}[label=\arabic*)]
    \item \textbf{Variables connues et recherchées :} $L = \qty{1}{mm} = \qty{1e-3}{m}$, $\lambda = \qty{\v{lam}e-9}{m}$ ; variable recherchée : $N$
    \item \textbf{Équations pertinentes :}
    \[ N = \frac{L}{\lambda} \]
    \item \textbf{Isolation de la variable recherchée :} l'équation est déjà isolée.
    \item \textbf{Résolution numérique :}
    \[ N = \frac{\qty{1e-3}{m}}{\qty{\v{lam}e-9}{m}} = \num[round-mode=figures,round-precision=3]{\fpeval{1e-3/\vsi{lam}}} \]
    \item \textbf{Réponse :} un millimètre contient environ \textbf{\num[round-mode=figures,round-precision=3]{\fpeval{1e-3/\vsi{lam}}}} longueurs d'onde.
\end{enumerate}
